\begin{lemma}[Situation 7.3, Case 1: component-intersection shortening]
Fix a finite loopless multigraph with exact edge identities and the fixed
rung-local minimum hierarchy-counterexample data.  Let \(C\) be the
original counterexample, let \(P\) be a
\noderef{HierarchyReplacementPresentation}, and write
\[
 \mu=P.\mathit{coloring},\quad
 H=P.\mathit{counterexample.hierarchy},\quad
 g=P.\mathit{counterexample.good},\quad
 q=H.\mathit{dividerCount},
\]
\[
 j=\noderef{HierarchyPathNumber}(C.\mathit{hierarchy}),\quad
 \ell=\noderef{HierarchyFullTailLength}(C),\quad
 A=\noderef{HierarchyStarredStage}(H,q),
\]
\[
 B=\noderef{HierarchyFullTailPrefix}(C,j-2),\qquad
 X=\noderef{HierarchyFullTailPrefix}(C,j-1).
\]
Let \(T_n\) be the last ladder tree of \(P\)'s replacement ETT, and put
\[
 D=D_{n,q}=\noderef{HierarchyStoredColorsAt}(H,q),\qquad
 \Gamma=\bigcup_{\eta\in D}g.\Gamma^q_\eta,
\]
\[
 M_0=\bigcup_{v\in
   V(P.\mathit{replacementBase.core})\setminus V(T_n)}
        \overline\mu(v).
\]
Define the finite set of colors on the exact displayed tail through
index \(j\) by
\[
 U_j=\{c:\text{ for some }r<|T|,\ |A|\le r<|A|+j
                 \text{ and }\mu(T.\mathit{edge}(r))=c\}.
\]
Thus each quantified position \(r\) selects the individual edge identity
\(T.\mathit{edge}(r)\), including when parallel edges have the same
endpoints.

Assume \(2\le j<\ell\).  Let \(\alpha\) and \(\beta\) be colors such that
\[
 \alpha\in\overline\mu\bigl(T.\mathit{vertex}(\operatorname{last})\bigr),
 \qquad \alpha\in\overline\mu(B),
\]
and such that
\[
 \alpha\notin\Gamma\cup M_0,
 \quad\hbox{or}\quad
 \alpha\in g.\Gamma^q_\eta\text{ for some }\eta\in D
 \text{ whose first \(\mu\)-missing position in \(T\) is at most }|B|.
\]
Suppose also that
\[
 \beta\ne\alpha,\qquad
 \beta\in\overline\mu(B),\qquad
 \beta\notin\Gamma,\qquad
 \beta\notin U_j,
\]
and that
\[
 \beta\notin M_0\cup D,
 \quad\hbox{or}\quad
 \beta\in D\ \text{ and }
   (g.\Gamma^q_\beta\cup\{\beta\})\cap U_j=\varnothing.
\]
Let
\[
 Q=\noderef{TwoColorComponent}(G,\mu,\alpha,\beta,y_\ell).
\]
If \(Q\) contains a vertex of \(X=T(y_{j-1})\), then a contradiction
follows from the lexicographic minimality of the current presentation.
\end{lemma}
\begin{proof}
Unfolding the data package
\noderef{HierarchyReplacementPresentation}, all the replacement objects
used below---the coloring \(\mu\), corresponding ETT, maximum-property
witness, base, counterexample, same-hierarchy certificate, and minimum
context---belong to one presentation on the same exact terminal tree \(T\).
By \noderef{HierarchyStarredStage}, \(A\) is the last starred stage of that
replacement hierarchy.  The two inequalities imply that \(B\) and \(X\)
are the exact prefixes of
\(T\) with respective lengths \(|A|+j-2\) and \(|A|+j-1\).  This follows
directly from \noderef{HierarchyFullTailLength},
\noderef{HierarchyFullTailPrefix}, and \(j<\ell=|T|-|A|\); in particular
no truncation by the minimum in the prefix definition occurs.  Hence
\[
 B\prec X\prec T,\qquad |X|=|B|+1,                           \tag{1}
\]
and the unique new vertex of \(X-B\) is \(y_{j-1}\).  The equal-stage
part of \(P.\mathit{sameGood}\), which is a
\noderef{SameGoodHierarchy}, identifies \(A\) and all exact prefixes here
with their counterparts in the replacement hierarchy.

Set \(y=y_\ell=T.\mathit{vertex}(\operatorname{last})\).  Because
\(j-1<\ell\), vertex injectivity in \noderef{TreeSequence} gives
\(y\notin V(X)\).  By \noderef{TreeMissingColors}, membership of
\(\alpha\) and \(\beta\) in \(\overline\mu(B)\) supplies vertices of
\(B\) at which they are respectively missing.  The exact prefix inclusion
\(B\preccurlyeq X\) retains those vertices, so the same definition shows
that both colors are missing on \(X\).  The coloring \(\mu\) is proper by the
\noderef{HierarchyMinimumCounterexampleContext} carried by \(P\).
Apply \noderef{TwoColorComponentTashkinovSplice} to \(X\), the root
\(y\), and the assumed intersection.  It chooses the first vertex \(u\)
of \(Q\) encountered from \(y\), and gives the reversed exact component
segment from \(u\) to \(y\).

If \(u=y_{j-1}\), append that segment to \(X\).  Otherwise, (1) and the
fact that no later vertex of the segment lies in \(X\) show that
\(u\in V(B)\).  Apply the splice lemma again with base \(B\), and append
its reversed first-intersection segment there.  Denote the resulting
tree-sequence in either case by \(T'\), and write \(S=X\) in the first case
and \(S=B\) in the second.  Thus
\[
 \begin{cases}
 X\preccurlyeq T'&\text{and the first new edge attaches at }y_{j-1},
       \quad u=y_{j-1},\\
 B\preccurlyeq T'&\text{and only the first new edge may attach at a
       nonterminal vertex of }B,\quad u\ne y_{j-1}.
 \end{cases}                                                  \tag{2}
\]
In both cases the appended exact edges have color \(\alpha\) or \(\beta\),
the appended vertices are new, the last vertex is \(y\), and the splice
lemma supplies a \(\mu\)-Tashkinov augmentation sequence from the base in
(2) to \(T'\).

We next verify the only exceptional seed issue before invoking the
mod-coloring theorem.  Let \(\mathcal E=P.\mathit{corresponding}.
\mathit{replacementETT}\), and let \(T_n\) be its terminal ladder tree.
Suppose simultaneously that \(q=0\), \(j=2\), and the last extension kind
is SE or RE.  At divider zero \noderef{HierarchyStarredStage} identifies
\(A\) with the core of \(P.\mathit{replacementBase}\), and the non-PE
clause of \noderef{HierarchyBase} identifies that core with \(T_n\).
The exact-prefix formula for \(B\) at \(j=2\) then gives \(A=B=T_n\),
while the eligible seed in
\noderef{ModColoring} is the exact one-edge sequence \(T_n+f_n=X\).
We claim that \(Q\) cannot meet \(T_n\).

Indeed, \(\alpha\) and \(\beta\) are both missing on \(T_n=B\).
Before applying the change lemma, derive fact (6.6) for the replacement
presentation: feed the properness, exact-one-uncoloredness, lower-rung,
fewer-divider, lexicographic-minimality, and large-missing-set fields of
\(P.\mathit{minimumContext}\) to
\noderef{HierarchyCounterexamplePenultimateElementary}.  The equal-stage
field of \noderef{SameGoodHierarchy} identifies its exact penultimate
full-tail prefix with the prefix used here.

Let \(\rho\) be the coloring immediately before the last series step.
The preceding-closed field of \noderef{TashkinovSeriesStep} says that
\(T_n\) is closed under \(\rho\).  In RE the current-coloring equation is
\(\mu=\rho\).  In SE, the current coloring \(\mu\) is the selected
coloring in a \noderef{MaximumDefectiveChoice}, whose selected-stable
field makes it stable from \(\rho\) on \(T_n\); both colorings are proper
and leave exactly \(e_0\) uncolored.  The vertex clause of stability
identifies their missing sets on \(T_n\).  If a boundary edge had color
\(\alpha\) under \(\mu\), then the reverse implication in
\noderef{sc1}, applied because \(\alpha\) is missing on \(T_n\), would
force that exact edge to have color \(\alpha\) under \(\rho\), contrary
to preceding closedness.  Thus in either SE/RE case
\noderef{ColoredVertexSetClosed} implies that \(T_n\) is
\(\alpha\)-closed in the singleton sense of
\noderef{TreeColorSetClosed}.  The ordered containment of \(T_n=A\) in
the exact penultimate full-tail prefix retains the two missing witnesses,
so \noderef{TreeMissingColors} gives the two missing-color premises of the
change lemma there.  Its conditional order-and-unused premise has false
antecedent because both colors are already missing on \(A\).  Apply the
closed-starred-stage clause of \noderef{ChangeLemma} to
\(P.\mathit{counterexample}\) and the six fields of
\(P.\mathit{minimumContext}\).  It says that \(T_n\) is
\((\alpha,\beta)\)-interchangeable and identifies the components through
the first \(\alpha\)- and \(\beta\)-missing witnesses in \(T_n\).

Call those first-missing witnesses \(a\) and \(b\), and denote their full
two-color components by \(K_a\) and \(K_b\).  They are actual vertices of
\(T_n\): the two colors are tree-missing there, and the least positions in
\noderef{HierarchyFirstMissingPosition} occur no later than the supplied
missing witnesses.  The root vertex \(y\) belongs to \(Q\).

We verify the path-component premises before using interchangeability.
The inequality \(\alpha\ne\beta\), the fact that \(a\in K_a\), and
\noderef{MissingColorIncidenceEndpoint} at the \(\alpha\)-missing vertex
\(a\) show that the neighbor set of \(a\) in the two-color incidence graph
has at most one element.  Thus the definition
\noderef{TwoColorPathComponent} makes \(K_a\) a path component.  Apply the
same endpoint lemma with the colors in the opposite order at the
\(\beta\)-missing vertex \(b\).  The adjacency formula in
\noderef{TwoColorIncidenceGraph} is unchanged when \(\alpha,\beta\) are
exchanged, so the resulting degree bound is a bound in the original graph;
therefore \(K_b\) is also a \(\noderef{TwoColorPathComponent}\).  Finally,
\(\alpha\) is missing at \(y\), so the endpoint lemma gives degree at most
one at \(y\).  Since \(y\in Q\), the same path-component definition makes
\(Q\) a two-color path component as well.

If \(a=b\), both colors are missing at this common vertex.  Were a node
adjacent to it, \noderef{TwoColorIncidenceGraph} would make that node an
exact incident edge colored \(\alpha\) or \(\beta\); either alternative is
in \noderef{PresentColors} at the common vertex and is therefore excluded
by the complementary definition \noderef{MissingColors}.  Thus its neighbor
set is empty, and \noderef{TwoColorComponent} makes \(K_a=K_b\) the isolated
one-vertex component.  If \(Q\) met \(T_n\), its certified path-component
status and that of \(K_a\), together with their respective intersections
with \(T_n\), would let \noderef{TreeInterchangeable} identify \(Q=K_a\).
This is impossible because \(Q\) contains the distinct vertex \(y\).

If \(a\ne b\), the preceding applications of
\noderef{MissingColorIncidenceEndpoint} give incidence degree at most one
at each of \(a,b,y\).  Were \(Q\) to meet \(T_n\), its path-component
status and that of \(K_a\) satisfy the two literal path premises of
\noderef{TreeInterchangeable}; the assumed vertex of \(Q\cap V(T_n)\) and
the vertex \(a\in K_a\cap V(T_n)\) give its two literal tree-intersection
premises.  Hence interchangeability identifies \(Q=K_a\), while the
component equality supplied by
\noderef{ChangeLemma} gives \(K_a=K_b\).  The three distinct nodes
\(a,b,y\) would then lie in one component and all have degree at most one.
But
\noderef{TwoColorComponentSpanningPath} spans that component by a simple
path.  Injectivity of the path order leaves only its two end nodes with
degree at most one: every internal node has its two distinct path
neighbors.  This is a contradiction.  Thus \(Q\cap V(T_n)=\varnothing\).

Since \(Q\) does meet \(X=T_n+f_n\), its first intersection with \(X\)
is the new vertex \(y_1=y_{j-1}\).  Hence only the first line of (2) can
occur in this exceptional case, and the exact edge identity \(f_n\) is
retained in \(T'\).  In every other case the appropriate source lies
below the selected splice base: for PE it is \(T_n\); for SE or RE with
positive divider count, \noderef{HierarchyNonparallelFirstEdge} puts the
exact edge \(f_n\) in the first hierarchy stage and hence below \(A\);
and when \(q=0\) but \(j>2\), the exact seed \(T_n+f_n\) already lies
below \(B\).                                                     \tag{3}

We now obtain an actual ETT on \(T'\).  The coloring \(\mu\) is stable
with itself on \(T_n\), directly from the two clauses of
\noderef{StableColoring}.  Apply \noderef{Td} to \(\mathcal E\), using
\(P.\mathit{maximum}\), positivity from
\(P.\mathit{counterexample}\), and the full lower-rung package from
\(P.\mathit{minimumContext}\).  We obtain
\[
                         \noderef{ModColoring}(\mathcal E,\mu).\tag{4}
\]
Inspect the last certified \noderef{TashkinovSeriesStep} of \(\mathcal E\).
In PE it contains a \(\mu\)-augmentation sequence from \(T_n\) to its
last closure.  In SE/RE it contains the singleton equality
\(F_n=\{f_n\}\), the exact \noderef{TreeSequenceAdjoinEdge}
\(T_n+f_n\), and the augmentation sequence from that seed.  The terminal
tree \(T\), hence each prefix in (2), is an exact initial segment of the
last closure by the \(\mathit{belowFinal}\) field in the definition
\noderef{ExtendedTashkinovTree}.  Apply
\noderef{TashkinovAugmentationSequenceRestrict}, using (3), to restrict
the certified route to the selected splice base.  Concatenate that route
with the augmentation sequence supplied by the splice lemma; this is
transitivity of the reflexive-transitive closure in
\noderef{TashkinovAugmentationSequence}.  The result is precisely the PE
or exact SE/RE eligibility route required by (4).

For its separate strict-extension premise, hierarchy nesting and the
PE/non-PE clauses of \noderef{HierarchyBase} give
\(T_n\preccurlyeq A\), while \(j\ge2\), the exact prefix formulas, and the
definition of \(S\) give \(A\preccurlyeq B\preccurlyeq S\).  Hence
\[
                         |T_n|\le |S|.
\]
The splice is strict over its selected base, so \(|S|<|T'|\).  Therefore
\[
                         |T_n|\le |S|<|T'|,
\]
and in particular \(|T_n|<|T'|\), exactly the strictness required by
\noderef{ModColoring}.  This argument also covers \(j=2\); the exceptional
SE/RE discussion above is needed only to retain the exact seed edge.

The universal clause of (4) therefore supplies an actual
\noderef{CorrespondingExtendedTashkinovTree} \(\mathcal E'\) with
terminal tree exactly \(T'\) and generating coloring \(\mu\).  Apply
\noderef{MaximumPropertyTransport} to this certificate and
\(P.\mathit{maximum}\); then
\[
                 \mathcal E'\text{ has maximum property}.    \tag{5}
\]
Notice that no replacement presentation, and no minimum context for
\(T'\), has been constructed.

We construct the base and hierarchy needed to compare \(T'\) with the
current minimum.  Correspondence identifies the rung number, each ladder
tree through \(T_n\), every raw stored set, every exact connecting-edge
set, and every extension kind.  By the definition
\noderef{HierarchyStoredColorsAt}, equality of those raw stored sets, the
ladder tree, and the vertexwise missing colors also identifies every stored
set \(D_{n,r}\) used by the copied hierarchy.  In the non-PE cases copy
\(P.\mathit{replacementBase}\) across these equalities; its core is
\(T_n\), and its nonparallel clause is unchanged.  In PE, the
corresponding certificate need not literally copy the internal chosen
vertex.  Extract the old and new last PE constructors from
\noderef{TashkinovSeriesStep}.  Connecting-edge-set equality identifies
their individual chosen edges \(f_n\).  Both maximum-defective records
write the endpoints of that same edge as an inside vertex of the common
exact \(T_n\) and an outside vertex.  Equality of the two unordered
endpoint pairs cannot exchange inside and outside, so the inside chosen
vertices agree.  Thus the old support position still displays the new
chosen vertex.  The old closure under the unchanged coloring \(\mu\) and
the old \noderef{OrderedTreeUnion} are unchanged.  This verifies every
field of a new \noderef{HierarchyBase} for \(\mathcal E'\), including the
terminal-ladder missing-set equality.

Keep the old divider count and every nonterminal stage and divider
position of \(H\), and replace only its terminal stage by \(T'\).  Every
old stage lies below \(A\), while (2) puts \(A\) strictly below \(T'\).
The nesting, strictness, prefix, and core-containment fields of
\noderef{TashkinovHierarchy} follow, with individual edge identities
unchanged.  Call the resulting hierarchy \(H'\).

Use exactly the same Gamma pairs as \(g\), indexed by the common stored
sets just computed from \noderef{HierarchyStoredColorsAt}.  The pair-size,
pair-disjointness, unique-refresh, PE-base, and internal minus-closed
fields of \noderef{GoodHierarchy} mention only the copied ladder tree,
base, nonterminal stages, stored sets, and Gamma pairs, so they transfer
literally.  We check the sole potentially changed clause, terminal
segment-color avoidance.  By the splice lemma every edge in \(T'\) beyond
the selected old prefix has color \(\alpha\) or \(\beta\).  The color
\(\beta\) belongs to no last Gamma pair by its common-choice certificate.
If \(\alpha\) satisfies the outside alternative, it too belongs to no
last Gamma pair.  Otherwise \(\alpha\in g.\Gamma^q_\eta\) for a stored
\(\eta\) whose first missing position is at most \(|B|\).  Since
\(|B|<|T|\), that position is an actual missing witness, not the fallback
last position in \noderef{HierarchyFirstMissingPosition}.  The same
witness lies in \(T'\), so the cutoff in
\noderef{HierarchySegmentColors} occurs before every appended edge.
Pairwise disjointness excludes \(\alpha\) from all other Gamma pairs.

For completeness, compare the old and new segment-color sets.  On the
common prefix through \(B\), the ordered edges and coloring values are
identical.  If an owner's first missing witness lies in that prefix, the
least such position is identical in the two trees.  Otherwise the new
first missing position either occurs on the appended segment, which can
only shorten the relevant cutoff after \(B\), or it is the new fallback
last position.  In all cases every new segment color is either an old
segment color or one of \(\alpha,\beta\).  The preceding paragraph
excludes precisely those two possible additions from the relevant Gamma
pair.  Hence the old avoidance field proves the new one.  We have obtained
a genuine good hierarchy on \(T'\), not merely a correspondence assertion.

The tree \(T'\) is not elementary.  One \(\alpha\)-missing witness lies
in \(B\), while \(y\) is also \(\alpha\)-missing and is the last vertex
of \(T'\).  They are distinct by (1) and vertex injectivity.  Unfolding
\noderef{ColoredVertexSetElementary}, the common color \(\alpha\) violates
disjointness of their two missing-color sets.  Together with positivity,
(5), and the good hierarchy just built, this defines a
\noderef{HierarchyCounterexample} \(C'\) on \(T'\).

Finally compare path numbers.  If the first line of (2) was used, every
edge from displayed position \(|A|+j-1\) onward joins consecutive
vertices, by the endpoint equations of the splice lemma.  If the second
line was used, its possibly nonconsecutive first appended edge has position
\(|B|=|A|+j-2\); every edge from the next position
\(|A|+j-1\) onward again joins consecutive vertices.  The strict splice
ensures that \(j-1\) belongs to the candidate interval in the definition
\noderef{HierarchyPathNumber}.  Hence
\[
 \noderef{HierarchyPathNumber}(H')\le j-1<j
       =\noderef{HierarchyPathNumber}(C.\mathit{hierarchy}). \tag{6}
\]
The corresponding certificate gives equality of rung numbers, and \(H'\)
has the same divider count as \(H\).  The
\noderef{SameGoodHierarchy} field of \(P\) identifies the original and
replacement divider counts, starred stages, and exact terminal tree.
Consequently their \noderef{HierarchyPathNumber}s and terminal excesses
are equal, so
\[
 \noderef{HierarchyCounterexampleMeasure}(P.\mathit{counterexample})
   =\noderef{HierarchyCounterexampleMeasure}(C).              \tag{7}
\]
Equations (6)--(7) make
\[
 \noderef{HierarchyCounterexampleLess}
   (\noderef{HierarchyCounterexampleMeasure}(C'),
    \noderef{HierarchyCounterexampleMeasure}
       (P.\mathit{counterexample}))
\]
true in its first coordinate.  This contradicts the
\(\mathit{lexicographicMinimality}\) field of the current
\noderef{HierarchyMinimumCounterexampleContext}
\(P.\mathit{minimumContext}\).  The contradiction proves the lemma.
\end{proof}
